% Appendices A--C (included after \appendix in thesis.tex).

\chapter{Extended tables for macro and micro metrics}
\label{app:tables}

The following tables consolidate every model--dataset cell written by the evaluation pipeline under \texttt{Thesis\_Research/outputs/results/}. Values are copied verbatim from the CSV exports (rounded for display). Re-running \texttt{run\_all.py} on a different machine or seed may shift decimals slightly; the qualitative ordering of models is stable across repeated local runs.

% AUTO-GENERATED by Thesis_Research/scripts/export_appendix_tables.py

\section{Macro-level (catalog) metrics}
CSV sources under \path{Thesis_Research/outputs/results/}: \path{combined_macro_bias_metrics.csv} and \path{macro_bias_metrics_take-a-lot-dataset.csv} (merged as a \texttt{take-a-lot} row). RMSE appears only where the pipeline exports a rating predictor.

\begin{longtable}{@{}llrrrr@{}}
\caption{Macro metrics (part I): concentration and popularity alignment.}\label{tab:app-macro-1}\\
\toprule
\textbf{Dataset} & \textbf{Model} & \textbf{Gini} & \textbf{ARP} & \textbf{Cat.\ cov.} & \textbf{$\rho_{\mathrm{pop,rec}}$} \\
\midrule
\endfirsthead
\multicolumn{6}{@{}l@{}}{\tablename\ \thetable{} -- \textit{continued from previous page}}\\
\toprule
\textbf{Dataset} & \textbf{Model} & \textbf{Gini} & \textbf{ARP} & \textbf{Cat.\ cov.} & \textbf{$\rho_{\mathrm{pop,rec}}$} \\
\midrule
\endhead
\midrule
\multicolumn{6}{r@{}}{\textit{Continued on next page}}\\
\endfoot
\bottomrule
\endlastfoot
book-crossing & ALS & 0.9950 & 43.89 & 0.0111 & 0.1204 \\
book-crossing & LogReg & 0.4997 & 18.96 & 0.9312 & 0.2598 \\
book-crossing & RandomForest & 0.9678 & 16.54 & 0.2484 & -0.0217 \\
book-crossing & SVD & 0.9833 & 31.81 & 0.0901 & 0.1826 \\
book-crossing & TF-IDF & 0.6907 & 20.80 & 0.7458 & 0.2096 \\
book-crossing & k-NN & 0.9360 & 10.87 & 0.2722 & -0.2347 \\
lastfm-2k & ALS & 0.9928 & 130.63 & 0.0166 & 0.1857 \\
lastfm-2k & LogReg & 0.6829 & 47.15 & 0.9522 & 0.1717 \\
lastfm-2k & RandomForest & 0.8110 & 86.23 & 0.8188 & 0.1298 \\
lastfm-2k & SVD & 0.9712 & 138.93 & 0.1778 & 0.2449 \\
lastfm-2k & TF-IDF & 0.9013 & 132.32 & 0.4466 & 0.2212 \\
lastfm-2k & k-NN & 0.8741 & 12.52 & 0.4837 & -0.3730 \\
ml-100k & ALS & 0.9874 & 233.09 & 0.0330 & 0.2557 \\
ml-100k & LogReg & 0.5428 & 75.12 & 0.8872 & 0.1318 \\
ml-100k & RandomForest & 0.9448 & 146.91 & 0.2031 & 0.2920 \\
ml-100k & SVD & 0.9552 & 157.28 & 0.2092 & 0.4339 \\
ml-100k & TF-IDF & 0.9254 & 61.56 & 0.3220 & 0.0739 \\
ml-100k & k-NN & 0.9800 & 156.04 & 0.0851 & 0.3407 \\
ml-1m & ALS & 0.9958 & 1178.31 & 0.0187 & 0.2015 \\
ml-1m & LogReg & 0.6102 & 288.67 & 0.9304 & 0.1711 \\
ml-1m & RandomForest & 0.9829 & 673.02 & 0.0626 & 0.1899 \\
ml-1m & SVD & 0.9568 & 832.10 & 0.2564 & 0.5180 \\
ml-1m & TF-IDF & 0.9370 & 327.33 & 0.3620 & 0.2576 \\
ml-1m & k-NN & 0.9885 & 635.72 & 0.0709 & 0.2854 \\
take-a-lot & ALS & 0.9998 & 28.75 & 0.0003 & 0.0365 \\
take-a-lot & LogReg & 0.7332 & 2.22 & 0.5634 & 0.1564 \\
take-a-lot & RandomForest & 0.9952 & 2.77 & 0.0249 & 0.0250 \\
take-a-lot & SVD & 0.9818 & 11.52 & 0.1191 & 0.2908 \\
take-a-lot & TF-IDF & 0.7578 & 3.15 & 0.4896 & 0.2085 \\
take-a-lot & k-NN & 0.9628 & 5.67 & 0.1601 & 0.2236 \\
\end{longtable}


\begin{longtable}{@{}llrrrrr@{}}
\caption{Macro metrics (part II): tail mass, diversity, metadata entropy, RMSE.}\label{tab:app-macro-2}\\
\toprule
\textbf{Dataset} & \textbf{Model} & \textbf{Avg.\ pop.\ pct.} & \textbf{LT@20\%} & \textbf{Agg.\ div.} & \textbf{Tok.\ ent.} & \textbf{RMSE} \\
\midrule
\endfirsthead
\multicolumn{7}{@{}l@{}}{\tablename\ \thetable{} -- \textit{continued from previous page}}\\
\toprule
\textbf{Dataset} & \textbf{Model} & \textbf{Avg.\ pop.\ pct.} & \textbf{LT@20\%} & \textbf{Agg.\ div.} & \textbf{Tok.\ ent.} & \textbf{RMSE} \\
\midrule
\endhead
\midrule
\multicolumn{7}{r@{}}{\textit{Continued on next page}}\\
\endfoot
\bottomrule
\endlastfoot
book-crossing & ALS & 0.0979 & 0.8957 & 23 & 3.38 & --- \\
book-crossing & LogReg & 0.4450 & 0.2885 & 1923 & 7.85 & --- \\
book-crossing & RandomForest & 0.4887 & 0.2040 & 513 & 5.63 & --- \\
book-crossing & SVD & 0.2371 & 0.5696 & 186 & 6.07 & 1.4963 \\
book-crossing & TF-IDF & 0.3876 & 0.3389 & 1540 & 6.93 & --- \\
book-crossing & k-NN & 0.6703 & 0.0458 & 562 & 7.17 & 1.7771 \\
lastfm-2k & ALS & 0.0601 & 0.9999 & 25 & 8.19 & --- \\
lastfm-2k & LogReg & 0.4603 & 0.2910 & 1435 & 9.16 & --- \\
lastfm-2k & RandomForest & 0.3651 & 0.3820 & 1234 & 9.38 & --- \\
lastfm-2k & SVD & 0.2407 & 0.5389 & 268 & 9.38 & 0.9060 \\
lastfm-2k & TF-IDF & 0.2100 & 0.6471 & 673 & 9.12 & --- \\
lastfm-2k & k-NN & 0.7216 & 0.0102 & 729 & 9.17 & 1.3924 \\
ml-100k & ALS & 0.0693 & 0.9389 & 38 & 3.31 & --- \\
ml-100k & LogReg & 0.4669 & 0.2310 & 1022 & 3.65 & --- \\
ml-100k & RandomForest & 0.2339 & 0.5345 & 234 & 3.39 & --- \\
ml-100k & SVD & 0.1809 & 0.5903 & 241 & 3.38 & 0.9311 \\
ml-100k & TF-IDF & 0.5091 & 0.1556 & 371 & 2.71 & --- \\
ml-100k & k-NN & 0.1651 & 0.6420 & 98 & 3.19 & 1.0041 \\
ml-1m & ALS & 0.0432 & 0.9993 & 61 & 2.73 & --- \\
ml-1m & LogReg & 0.4518 & 0.2534 & 3033 & 3.78 & --- \\
ml-1m & RandomForest & 0.2582 & 0.5327 & 204 & 3.53 & --- \\
ml-1m & SVD & 0.1923 & 0.6882 & 836 & 3.46 & 0.8726 \\
ml-1m & TF-IDF & 0.4058 & 0.2869 & 1180 & 2.86 & --- \\
ml-1m & k-NN & 0.2244 & 0.6156 & 231 & 3.53 & 0.9769 \\
take-a-lot & ALS & 0.0026 & 1.0000 & 14 & 6.09 & --- \\
take-a-lot & LogReg & 0.4633 & 0.2028 & 27556 & 10.01 & --- \\
take-a-lot & RandomForest & 0.4110 & 0.2890 & 1217 & 8.44 & --- \\
take-a-lot & SVD & 0.1015 & 0.8215 & 5824 & 9.34 & 0.9662 \\
take-a-lot & TF-IDF & 0.3940 & 0.3045 & 23945 & 10.07 & --- \\
take-a-lot & k-NN & 0.2251 & 0.5749 & 7830 & 9.61 & 1.0211 \\
\end{longtable}


\section{Micro-level (user-centric) metrics}
Segment metrics mirror \path{combined_user_centric_metrics.csv} and \path{user_centric_metrics_take-a-lot-dataset.csv}. $\Delta$RMSE is niche minus mainstream RMSE (positive implies higher error on niche users).

\begin{longtable}{@{}llrrrr@{}}
\caption{Micro metrics (part I): RMSE by taste segment and $\Delta$RMSE.}\label{tab:app-micro-1}\\
\toprule
\textbf{Dataset} & \textbf{Model} & \textbf{RMSE\textsubscript{niche}} & \textbf{RMSE\textsubscript{div.}} & \textbf{RMSE\textsubscript{main}} & \textbf{$\Delta$RMSE} \\
\midrule
\endfirsthead
\multicolumn{6}{@{}l@{}}{\tablename\ \thetable{} -- \textit{continued from previous page}}\\
\toprule
\textbf{Dataset} & \textbf{Model} & \textbf{RMSE\textsubscript{niche}} & \textbf{RMSE\textsubscript{div.}} & \textbf{RMSE\textsubscript{main}} & \textbf{$\Delta$RMSE} \\
\midrule
\endhead
\midrule
\multicolumn{6}{r@{}}{\textit{Continued on next page}}\\
\endfoot
\bottomrule
\endlastfoot
book-crossing & ALS & 1.7427 & 1.6719 & 1.7295 & 0.0132 \\
book-crossing & LogReg & 1.7427 & 1.6719 & 1.7295 & 0.0132 \\
book-crossing & RandomForest & 1.5814 & 1.5817 & 1.6347 & -0.0533 \\
book-crossing & SVD & 1.4845 & 1.4927 & 1.5165 & -0.0320 \\
book-crossing & TF-IDF & 1.7427 & 1.6719 & 1.7295 & 0.0132 \\
book-crossing & k-NN & 1.8546 & 1.7226 & 1.7419 & 0.1127 \\
lastfm-2k & ALS & 1.3925 & 1.4177 & 1.4172 & -0.0247 \\
lastfm-2k & LogReg & 1.3925 & 1.4177 & 1.4172 & -0.0247 \\
lastfm-2k & RandomForest & 1.1766 & 1.2507 & 1.2454 & -0.0688 \\
lastfm-2k & SVD & 0.8364 & 0.9051 & 0.9503 & -0.1140 \\
lastfm-2k & TF-IDF & 1.3925 & 1.4177 & 1.4172 & -0.0247 \\
lastfm-2k & k-NN & 1.3947 & 1.4217 & 1.3643 & 0.0304 \\
ml-100k & ALS & 1.1354 & 1.0948 & 1.0721 & 0.0633 \\
ml-100k & LogReg & 1.1354 & 1.0948 & 1.0721 & 0.0633 \\
ml-100k & RandomForest & 1.0442 & 1.0221 & 1.0201 & 0.0241 \\
ml-100k & SVD & 0.9346 & 0.9253 & 0.9294 & 0.0052 \\
ml-100k & TF-IDF & 1.1354 & 1.0948 & 1.0721 & 0.0633 \\
ml-100k & k-NN & 1.0235 & 0.9802 & 0.9755 & 0.0480 \\
ml-1m & ALS & 1.1471 & 1.0800 & 1.0748 & 0.0723 \\
ml-1m & LogReg & 1.1471 & 1.0800 & 1.0748 & 0.0723 \\
ml-1m & RandomForest & 1.0191 & 0.9696 & 0.9671 & 0.0520 \\
ml-1m & SVD & 0.8740 & 0.8677 & 0.8772 & -0.0032 \\
ml-1m & TF-IDF & 1.1471 & 1.0800 & 1.0748 & 0.0723 \\
ml-1m & k-NN & 0.9975 & 0.9482 & 0.9493 & 0.0482 \\
take-a-lot & ALS & 1.0235 & 1.0655 & 0.9467 & 0.0769 \\
take-a-lot & LogReg & 1.0235 & 1.0655 & 0.9467 & 0.0769 \\
take-a-lot & RandomForest & 1.0229 & 1.0610 & 0.9409 & 0.0819 \\
take-a-lot & SVD & 0.9663 & 1.0218 & 0.9059 & 0.0604 \\
take-a-lot & TF-IDF & 1.0235 & 1.0655 & 0.9467 & 0.0769 \\
take-a-lot & k-NN & 1.0249 & 1.0763 & 0.9572 & 0.0677 \\
\end{longtable}


\begin{longtable}{@{}llrrr@{}}
\caption{Micro metrics (part II): NDCG@10 by taste segment.}\label{tab:app-micro-2}\\
\toprule
\textbf{Dataset} & \textbf{Model} & \textbf{NDCG\textsubscript{niche}} & \textbf{NDCG\textsubscript{div.}} & \textbf{NDCG\textsubscript{main}} \\
\midrule
\endfirsthead
\multicolumn{5}{@{}l@{}}{\tablename\ \thetable{} -- \textit{continued from previous page}}\\
\toprule
\textbf{Dataset} & \textbf{Model} & \textbf{NDCG\textsubscript{niche}} & \textbf{NDCG\textsubscript{div.}} & \textbf{NDCG\textsubscript{main}} \\
\midrule
\endhead
\midrule
\multicolumn{5}{r@{}}{\textit{Continued on next page}}\\
\endfoot
\bottomrule
\endlastfoot
book-crossing & ALS & 0.0124 & 0.0123 & 0.0116 \\
book-crossing & LogReg & 0.0592 & 0.0502 & 0.0408 \\
book-crossing & RandomForest & 0.0095 & 0.0076 & 0.0085 \\
book-crossing & SVD & 0.0024 & 0.0085 & 0.0101 \\
book-crossing & TF-IDF & 0.1120 & 0.0821 & 0.0737 \\
book-crossing & k-NN & 0.0008 & 0.0016 & 0.0029 \\
lastfm-2k & ALS & 0.0210 & 0.0509 & 0.0254 \\
lastfm-2k & LogReg & 0.0310 & 0.0284 & 0.0743 \\
lastfm-2k & RandomForest & 0.0132 & 0.0155 & 0.0642 \\
lastfm-2k & SVD & 0.0157 & 0.0215 & 0.0914 \\
lastfm-2k & TF-IDF & 0.1507 & 0.1403 & 0.2557 \\
lastfm-2k & k-NN & 0.0000 & 0.0008 & 0.0001 \\
ml-100k & ALS & 0.1709 & 0.1004 & 0.0727 \\
ml-100k & LogReg & 0.0653 & 0.0305 & 0.0151 \\
ml-100k & RandomForest & 0.1044 & 0.0644 & 0.0363 \\
ml-100k & SVD & 0.1465 & 0.0805 & 0.0435 \\
ml-100k & TF-IDF & 0.0660 & 0.0346 & 0.0234 \\
ml-100k & k-NN & 0.1514 & 0.0828 & 0.0416 \\
ml-1m & ALS & 0.1396 & 0.0902 & 0.0508 \\
ml-1m & LogReg & 0.0396 & 0.0193 & 0.0120 \\
ml-1m & RandomForest & 0.0604 & 0.0395 & 0.0251 \\
ml-1m & SVD & 0.1274 & 0.0875 & 0.0572 \\
ml-1m & TF-IDF & 0.0697 & 0.0549 & 0.0327 \\
ml-1m & k-NN & 0.0792 & 0.0434 & 0.0231 \\
take-a-lot & ALS & 0.0002 & 0.0010 & 0.0017 \\
take-a-lot & LogReg & 0.0038 & 0.0064 & 0.0043 \\
take-a-lot & RandomForest & 0.0000 & 0.0000 & 0.0004 \\
take-a-lot & SVD & 0.0001 & 0.0001 & 0.0015 \\
take-a-lot & TF-IDF & 0.0117 & 0.0118 & 0.0175 \\
take-a-lot & k-NN & 0.0013 & 0.0011 & 0.0004 \\
\end{longtable}


\section{Longitudinal feedback simulation (full iteration grid)}
All iterations for SVD and TF--IDF from \path{combined_feedback_simulation_metrics.csv}. The main text stresses qualitative trajectories; this grid supports exact reproduction of line charts.

\begin{longtable}{@{}llrrr@{}}
\caption{Feedback simulation: aggregate diversity and ARP at every iteration (SVD vs.\ TF--IDF).}\label{tab:app-feedback-full}\\
\toprule
\textbf{Dataset} & \textbf{Model} & \textbf{Iter.} & \textbf{Agg.\ diversity} & \textbf{ARP} \\
\midrule
\endfirsthead
\multicolumn{5}{@{}l@{}}{\tablename\ \thetable{} -- \textit{continued from previous page}}\\
\toprule
\textbf{Dataset} & \textbf{Model} & \textbf{Iter.} & \textbf{Agg.\ diversity} & \textbf{ARP} \\
\midrule
\endhead
\midrule
\multicolumn{5}{r@{}}{\textit{Continued on next page}}\\
\endfoot
\bottomrule
\endlastfoot
book-crossing & SVD & 1 & 134 & 32.79 \\
book-crossing & SVD & 2 & 306 & 26.26 \\
book-crossing & SVD & 3 & 533 & 24.34 \\
book-crossing & SVD & 4 & 609 & 24.50 \\
book-crossing & SVD & 5 & 715 & 23.63 \\
book-crossing & TF-IDF & 1 & 1264 & 21.29 \\
book-crossing & TF-IDF & 2 & 1330 & 24.78 \\
book-crossing & TF-IDF & 3 & 1348 & 27.52 \\
book-crossing & TF-IDF & 4 & 1382 & 30.13 \\
book-crossing & TF-IDF & 5 & 1404 & 32.16 \\
lastfm-2k & SVD & 1 & 170 & 176.66 \\
lastfm-2k & SVD & 2 & 103 & 245.06 \\
lastfm-2k & SVD & 3 & 129 & 560.07 \\
lastfm-2k & SVD & 4 & 148 & 661.76 \\
lastfm-2k & SVD & 5 & 166 & 699.74 \\
lastfm-2k & TF-IDF & 1 & 492 & 150.24 \\
lastfm-2k & TF-IDF & 2 & 516 & 174.06 \\
lastfm-2k & TF-IDF & 3 & 532 & 188.95 \\
lastfm-2k & TF-IDF & 4 & 531 & 197.58 \\
lastfm-2k & TF-IDF & 5 & 538 & 202.75 \\
ml-100k & SVD & 1 & 174 & 165.14 \\
ml-100k & SVD & 2 & 141 & 264.26 \\
ml-100k & SVD & 3 & 157 & 335.54 \\
ml-100k & SVD & 4 & 143 & 358.38 \\
ml-100k & SVD & 5 & 162 & 405.70 \\
ml-100k & TF-IDF & 1 & 235 & 68.34 \\
ml-100k & TF-IDF & 2 & 286 & 89.03 \\
ml-100k & TF-IDF & 3 & 333 & 101.94 \\
ml-100k & TF-IDF & 4 & 366 & 109.40 \\
ml-100k & TF-IDF & 5 & 394 & 115.63 \\
ml-1m & SVD & 1 & 649 & 893.35 \\
ml-1m & SVD & 2 & 541 & 918.07 \\
ml-1m & SVD & 3 & 523 & 1321.16 \\
ml-1m & SVD & 4 & 528 & 1832.50 \\
ml-1m & SVD & 5 & 529 & 1997.48 \\
ml-1m & TF-IDF & 1 & 805 & 348.38 \\
ml-1m & TF-IDF & 2 & 983 & 494.43 \\
ml-1m & TF-IDF & 3 & 1084 & 556.39 \\
ml-1m & TF-IDF & 4 & 1183 & 576.73 \\
ml-1m & TF-IDF & 5 & 1278 & 584.20 \\
\end{longtable}


\chapter{Additional visualizations and diagrams}
\label{app:viz}

This appendix collects schematic views of the evaluation workflow and reproduces the cross-dataset figures from the main experiments at appendix scale, so printed copies remain self-contained even when individual chapters are excerpted.

\section{End-to-end evaluation pipeline}
Figure~\ref{fig:app-pipeline} situates Experiments~A--C as three measurement stations on a shared data and training backbone: the same splits, models, and recommendation lists feed macro catalog metrics, segment-level user metrics, and the longitudinal click simulator.

\begin{figure}[p]
  \centering
  \begin{tikzpicture}[
    font=\small,
    box/.style={draw, rounded corners=2pt, align=center, inner sep=6pt, minimum width=3.1cm},
    arr/.style={-{Latex[length=2.2mm]}, thick}
  ]
    \node[box, fill=gray!12] (data) {Datasets\\\textit{ML-100K, ML-1M, Last.fm, Book-Crossing, Takealot}};
    \node[box, below=0.9cm of data] (prep) {Preprocess\\\textit{$k$-core, TF--IDF text, Surprise CF}};
    \node[box, below=0.9cm of prep] (train) {Train baselines\\\textit{ALS, SVD, $k$-NN, TF--IDF, RF, LogReg}};
    \node[box, below left=1.1cm and 0.2cm of train] (a) {\textbf{Exp.~A}\\Macro bias\\Gini, ARP, coverage};
    \node[box, below=1.1cm of train] (b) {\textbf{Exp.~B}\\Micro fairness\\RMSE / NDCG segments};
    \node[box, below right=1.1cm and 0.2cm of train] (c) {\textbf{Exp.~C}\\Feedback loop\\5 iterations};
    \draw[arr] (data) -- (prep);
    \draw[arr] (prep) -- (train);
    \draw[arr] (train.south west) -- (a.north);
    \draw[arr] (train.south) -- (b.north);
    \draw[arr] (train.south east) -- (c.north);
  \end{tikzpicture}
  \caption{Unified pipeline: shared preprocessing and model training feed three evaluation tracks (Chapters~\ref{ch:exp_a}--\ref{ch:exp_c}).}
  \label{fig:app-pipeline}
\end{figure}

\section{Stylised feedback loop with explicit simulator knobs}
Figure~\ref{fig:app-feedback-schematic} summarises the feedback simulation at the level of state variables referenced in Chapter~\ref{ch:exp_c}. Clicks use the linear blend $c_{ui}=\mathrm{clip}(0.7\, s_{ui}+0.3\, \tilde{p}_i,[0,1])$ from Section~\ref{sec:method_click}, then a Bernoulli draw on each recommended slot.

\begin{figure}[p]
  \centering
  \begin{tikzpicture}[
    font=\small,
    n/.style={draw, circle, minimum size=1.05cm, inner sep=0pt},
    sq/.style={draw, rounded corners=2pt, minimum width=2.6cm, minimum height=0.85cm, align=center},
    arr/.style={-{Latex[length=2mm]}, thick}
  ]
    \node[sq] (log) {$\mathcal{D}_{t-1}$\\training log};
    \node[sq, right=2.4cm of log] (rec) {Recommender\\top-$5$ lists};
    \node[sq, right=2.4cm of rec] (clk) {Clicks\\Bernoulli rule};
    \node[sq, below=1.5cm of rec] (logp) {$\mathcal{D}_t$\\updated log};
    \draw[arr] (log) -- node[above]{train / infer} (rec);
    \draw[arr] (rec) -- node[above]{expose} (clk);
    \draw[arr] (clk.south) -- ++(0,-0.35) -| node[pos=0.22, below]{append $5.0$ ratings} (logp.south);
    \draw[arr] (logp.west) -- ++(-0.6,0) |- node[pos=0.25, left]{next round} (log.west);
  \end{tikzpicture}
  \caption{Longitudinal simulator: logs evolve from recommendations and stochastic clicks, then re-enter training.}
  \label{fig:app-feedback-schematic}
\end{figure}

\section{Cross-dataset figures (same assets as the main text)}
For convenience, the heatmaps and segment plots used in Chapters~\ref{ch:exp_a}--\ref{ch:exp_c} are repeated here without reinterpretation.

\begin{figure}[p]
  \centering
  \includegraphics[width=\linewidth]{combined_macro_gini_heatmap.png}
  \caption[Appendix: Gini heatmap]{Gini coefficient of recommendation frequency versus item popularity (all non-Takealot benchmarks).}
\end{figure}

\begin{figure}[p]
  \centering
  \includegraphics[width=\linewidth]{combined_macro_catalog_coverage_heatmap.png}
  \caption[Appendix: catalog coverage heatmap]{Catalog coverage of top-$N$ recommendations across datasets and models.}
\end{figure}

\begin{figure}[p]
  \centering
  \includegraphics[width=0.92\textwidth]{combined_fairness_gap_delta_rmse.png}
  \caption[Appendix: niche RMSE gap]{Niche minus mainstream RMSE ($\Delta$RMSE) by dataset and model.}
\end{figure}

\begin{figure}[p]
  \centering
  \includegraphics[width=0.92\textwidth]{combined_rmse_niche_vs_mainstream.png}
  \caption[Appendix: RMSE by segment]{RMSE for niche versus mainstream taste segments.}
\end{figure}

\begin{figure}[p]
  \centering
  \includegraphics[width=0.92\textwidth]{combined_feedback_diversity_decay.png}
  \caption[Appendix: feedback diversity]{Aggregate diversity and ARP across five feedback iterations (SVD versus TF--IDF).}
\end{figure}

\chapter{Extended software specifications}
\label{appendix:software}

\section{Repository layout}
All experiments are implemented as a single Python package-style tree under \texttt{Thesis\_Research/}: \texttt{run\_all.py} orchestrates dataset ingestion, training, evaluation, plotting, and CSV exports; \texttt{src/} holds metric utilities and statistical tests; \texttt{outputs/results/} stores the tables consumed by this document. Figures are written to \texttt{outputs/figures/} (and mirrored under \texttt{Thesis-main/outputs/figures/} when syncing for \LaTeX).

\begin{figure}[htbp]
  \centering
  \begin{tikzpicture}[
    font=\footnotesize,
    every node/.style={draw, align=left, anchor=west, inner sep=4pt}
  ]
    \node (r) at (0,0) {\texttt{Thesis\_Research/}};
    \node (run) at (0.4,-0.85) {\texttt{run\_all.py}   entry point};
    \node (src) at (0.4,-1.55) {\texttt{src/}   models, metrics, \texttt{statistical\_tests.py}};
    \node (out) at (0.4,-2.25) {\texttt{outputs/results/*.csv}   thesis tables};
    \node (fig) at (0.4,-2.95) {\texttt{outputs/figures/*.png}   thesis figures};
    \draw[-{Latex}] (r.south west) ++(0.15,-0.12) |- (run.west);
    \draw[-{Latex}] (r.south west) ++(0.15,-0.12) |- (src.west);
    \draw[-{Latex}] (r.south west) ++(0.15,-0.12) |- (out.west);
    \draw[-{Latex}] (r.south west) ++(0.15,-0.12) |- (fig.west);
  \end{tikzpicture}
  \caption{Code and artefact layout relative to the research repository root.}
  \label{fig:app-repo-layout}
\end{figure}

\section{Language runtime and hardware assumptions}
The code targets \textbf{Python 3.10+} (tested on 64-bit Windows~10/11). Experiments are CPU-bound; GPU acceleration is not required. Wall-clock runtime scales with Book-Crossing and Takealot matrix sizes; a full multi-dataset sweep completes overnight on a modern laptop when Surprise and implicit builds are already compiled.

\section{Core dependency pins}
Table~\ref{tab:app-deps} lists the libraries that materially affect numerical results. A broader scientific stack (Jupyter, widgets, conda solver libraries) may be present in a full Anaconda install but does not change the reported metrics.

\begin{table}[htbp]
  \centering
  \small
  \begin{tabular}{@{}llp{6.8cm}@{}}
    \toprule
    \textbf{Package} & \textbf{Role in this thesis} & \textbf{Notes} \\
    \midrule
    \texttt{numpy} & Arrays, linear algebra & Pinned below 2.x for \texttt{implicit} compatibility. \\
    \texttt{pandas} & CSV aggregation & Used for all combined result tables. \\
    \texttt{scipy} & Statistics & Welch/Mann--Whitney helpers in \texttt{statistical\_tests.py}. \\
    \texttt{scikit-learn} & TF--IDF, RF, LogReg, metrics & Pipelines mirror standard sklearn idioms. \\
    \texttt{scikit-surprise} & SVD, $k$-NN, ALS & Optional import guarded when wheels are unavailable. \\
    \texttt{implicit} & ALS (implicit feedback path) & Used where the benchmark supports implicit-style training. \\
    \texttt{matplotlib}, \texttt{seaborn} & Plots & Heatmaps and line charts exported to PNG. \\
    \texttt{tqdm} & Progress bars & Long Takealot sweeps remain observable during batch runs. \\
    \bottomrule
  \end{tabular}
  \caption{Scientific stack directly exercised by \texttt{run\_all.py}.}
  \label{tab:app-deps}
\end{table}

\section{Pinned versions (\texttt{pip freeze} excerpt)}
Listing~\ref{lst:app-freeze} lists pinned versions from the development host. Reproduce with a clean virtual environment and the research repository \path{requirements.txt}, then align versions as needed.

\lstset{basicstyle=\ttfamily\footnotesize,frame=single,breaklines=true,columns=fullflexible}
\lstinputlisting[caption={Excerpted \texttt{pip} pins for core scientific dependencies.},label={lst:app-freeze}]{thesis_core_environment.txt}

\section{Regeneration checklist}
\begin{enumerate}
  \item Install dependencies from \path{requirements.txt} in the research repo; verify \texttt{surprise} imports.
  \item From the repo root, run \texttt{python run\_all.py} with the thesis defaults (top-$10$ lists for offline metrics; feedback uses top-$5$ and five iterations).
  \item Copy new PNGs into folder \path{Thesis-main/outputs/figures/} when compiling the thesis locally.
  \item Run \texttt{python scripts/export\_appendix\_tables.py} to rebuild \path{_appendix_generated_tables.tex}.
  \item Recompile \texttt{thesis.tex} (\texttt{pdflatex}, \texttt{bibtex}, \texttt{pdflatex} $\times 2$).
\end{enumerate}
